%!TEX program = uplatex
\RequirePackage{plautopatch}
\documentclass[uplatex,dvipdfmx]{jlreq}
\usepackage{amsmath,amssymb}

\pagestyle{empty}

\begin{document}

%======================================================================
% 表紙
%======================================================================
\begin{center}
  \vspace*{20pt}
  {\Huge\bfseries 曲面の写像類群}

  \vspace{10pt}

  {\Large\bfseries いろいろな数学の集いの広場}

  \vspace{10pt}

  {\large 概要}
\end{center}
\vfill

\newpage

%======================================================================
% 講演１：森田茂之「曲面の写像類群とは」
%======================================================================
\begin{center}
{\large\bfseries 曲面の写像類群とは}

\hspace*{\fill}{\large\bfseries 森田 茂之（東大・数理）}
\end{center}

\medskip

曲面、すなわち実２次元の多様体は現代幾何学の発祥の地といえよう。
実際、Gauss 曲率、Gauss--Bonnet の定理、Poincaré--Hopf の定理、
そしてまた種数による閉曲面の位相的な分類や Gauss 写像は、
微分幾何学、トポロジー、大域解析学などの分野に分かれて発展してきた現代幾何学の共通の原点である。

一方 Riemann により導入された Riemann 面は、
現在では１次元複素多様体と同義語となっているが、
幾何的には１変数の（多価）解析関数が真に定義されている場所として、
複素平面を切り貼りしてできる曲面のことである。
代数的には代数曲線、すなわち１次元代数多様体は、
代数幾何学、複素多様体論の故郷であり、さらには整数論とも深い関係がある。

高次元では上記の二つの立場、すなわち実多様体と複素（代数）多様体の世界は
かなり異なる様相を示すようになる。
しかし実２次元と複素１次元の世界はほとんど等価なのである。
このことはいわゆる１意化定理により、少なくとも種数が２以上の場合、
コンパクト Riemann 面＝非特異代数曲線＝負の定曲率曲面という等式が成り立つことに端的に示されている。

このように、さまざまな立場を超えた共通の存在である「曲面」の合同類をすべて集めてできる空間
（ただし種数は固定する）を、Riemann のモジュライ空間という。
この空間は非常に豊かな構造をもつことが、これまでの研究により次第に明らかになってきている。
しかし、まだまだ神秘的なヴェールをまとったままの部分の方がはるかに多い。

さて曲面の写像類群とは、その名の通り曲面の（向きを保つ）微分同相写像の
イソトピー類全体のなす群である。
このような群は任意の多様体に対して定義される。
しかし、曲面の場合は上記のモジュライ空間の基本群としての役割を果たすため、
際だって重要な研究対象となっている。

ここでは、この曲面の写像類群についてのできるだけ分かりやすい解説をしたい。

\newpage

%======================================================================
% 講演２：河澄響矢「森田・マンフォード類、その微分形式表示と『内部構造』」
%======================================================================
\begin{center}
{\large\bfseries 森田・マンフォード類、その微分形式表示と「内部構造」}

\hspace*{\fill}{\large\bfseries 河澄 響矢（北大・理）}
\end{center}

\medskip

写像類群のコホモロジー環は、（係数を有理数係数にする、つまり torsion を無視すると、）
リーマン面のモジュライ空間のコホモロジー環に一致している。
その最も重要な中身が森田茂之と D.~Mumford が独立に発見した
森田・マンフォード類、別称 tautological class である。
これは、モジュライ空間の上に乗っているリーマン面の普遍族の上で、
相対接束のオイラー類の冪をとり、ファイバー積分によりモジュライ空間に落としたものである。
この講演では、森田・マンフォード類をモジュライ空間の上で微分形式で表示する
二通りの方法を解説したい。

ひとつはリーマン面の普遍族に入っている双曲計量を使う方法で、
S.~Wolpert により実行された。
とくに、第一森田・マンフォード類がヴェイユ・ピーターソン・ケーラー形式によって表されることが分かる。

もう一つは森田・マンフォード類をねじれ係数に一般化することで見えてくる
「森田・マンフォード類の内部構造」を使う方法である。
第３種アーベル積分を写像類群の言葉で解釈すると、
２点付き曲面の写像類群のねじれ準同型が得られる。
これを相対接束のオイラー類と混ぜあわせて冪をとり、ファイバー積分すると、
森田・マンフォード類をねじれ係数に一般化することが出来る。
ここで「第３種アーベル積分」の $3$ 乗のファイバー積分は拡大ジョンソン準同型に一致する。
ねじれ係数森田・マンフォード類が係数の縮約で閉じていることから、
拡大ジョンソン準同型から得られる自明係数コホモロジー類の全体は
（もともとの）森田・マンフォード類の生成する多項式環に一致する。
したがって、第３種アーベル積分の第一変分の３乗のファイバー積分こそが、
拡大ジョンソン準同型の微分形式表示に他ならない。
これは、点付き調和体積の第一変分に一致する。
この微分形式を縮約することで森田・マンフォード類の
（双曲計量を使わない）微分形式表示が複数得られるのである。
これらの表示は（講演者がさぼっているために）分からないことだらけであるが、
繊細な情報を含んでおり、いろいろな幾何学との深い関係を予感させる。
双曲幾何学はリーマン面にとって唯一絶対の幾何学ではないようである。

\newpage

%======================================================================
% 講演３：森田茂之「写像類群をめぐるこれまでの結果と夢」
%======================================================================
\begin{center}
{\large\bfseries 写像類群をめぐるこれまでの結果と夢}

\hspace*{\fill}{\large\bfseries 森田 茂之（東大・数理）}
\end{center}

\medskip

写像類群が数学の多くの分野と関わる驚くほど豊富な構造をもっていることを示すのが、
今回の集会の主要な目的である。
ここでは主として写像類群のコホモロジーの観点から、
これまでに知られているいろいろな結果を概観し、
それらをもとにいくつかの予想と夢を述べてみたい。
主役を果たすのは、ナイーヴな意味で「写像類群の Lie 代数」ともいうべき、
ある無限次元の Lie 代数 $\ell_g$ である。
この無限次元 Lie 代数は、Kontsevich が無限次元シンプレクティック幾何学の
枠組みの中で扱った三つの Lie 代数のうちの一つであるが、
その構造はまだまだ未解明のままといってよい。
上記でナイーヴな意味でと書いたが、実はこの Lie 代数の中で写像類群
あるいは Riemann 面のモジュライ空間に対応する部分は比較的小さいことがわかってきている。
実際には、いわば「真の写像類群の Lie 代数」$\mathbf{m}_g$ から $\ell_g$ への
Johnson 準同型と呼ばれる射 $\tau : \mathbf{m}_g \to \ell_g$ があり、
その像がモジュライ空間に対応する部分となっている。

具体的にはつぎのようなテーマとそれらの間の関係について述べてみたい。
\begin{enumerate}
  \setlength{\itemsep}{2pt}
  \item モジュライ空間の安定特性類と Johnson 準同型の像
  \item モジュライ空間の非安定コホモロジーに関する Faber 予想
  \item Johnson 準同型の核と３次元多様体の位相不変量
  \item $\ell_g$ の生成元と関係式：Johnson 準同型と Trace
  \item $\ell_g$ のコホモロジーと自由群の外部自己同型群の特性類
  \item 写像類群の２次特性類とモジュライ空間の高次の幾何学
\end{enumerate}

\newpage

%======================================================================
% 講演４：阿原一志「写像類群とコンピュータ」
%======================================================================
\begin{center}
{\large\bfseries 写像類群とコンピュータ}

\hspace*{\fill}{\large\bfseries 阿原 一志（明治大・理工）}
\end{center}

\medskip

「猫も杓子もコンピュータ」時代をこれから迎える。
一方で写像類群の計算にもコンピュータが応用できて然るべきであると考えられて久しい。
解析などでは「コンピュータで計算することにより定理が証明できる」などという夢のような
次元に突入しているようだが、写像類群ではどうか。
この点について考えてみよう。問題点を次のように分類してみる。
\begin{enumerate}
  \setlength{\itemsep}{2pt}
  \item[(1)] 曲面の基本群や写像類群をコンピュータで扱うことができるか。
  \item[(2)] 写像類群の表現、コサイクル、コホモロジーを具体的に計算できるか。
  \item[(3)] 写像類群を勉強する人が楽しめるような視覚ソフトを作れるか。
\end{enumerate}

（1）について。これは単に無限離散群の問題を考えることで、
automatic 構造、語問題について考えることになるが、ここでは触れない。

（2）について。知られている表現やコサイクルについて、
その具体的な値を求めることにこしたことはない。
（難しい不変量だと具体的な値すらもわからない場合が多い。）
その結果を眺めることによって、何か法則性が知れるかも知れない。
講演では、この一例として、Meyer 関数、Jones 表現を計算する
Mathematica のプログラムを紹介する。
コンピュータで群のコホモロジーを求めるのは阿原にとって将来の課題である。

（3）について。写像類群自体を目で見ることは難しいが、
曲面上のカーブやデーンの捻れをコンピュータのグラフィック表示することはたやすい。
講演では、この一例として、「てるあき」というゲームソフトを紹介する。

\newpage

%======================================================================
% 講演５：村上順「ウェブ図からみた写像類群」
%======================================================================
\begin{center}
{\large\bfseries ウェブ図からみた写像類群}

\hspace*{\fill}{\large\bfseries 村上 順（阪大・理）}
\end{center}

\medskip

曲面の写像類群が豊富な構造を持つことを示すのがこの集会の主要な目的なのだが、
この講演では、ウェブ図を用いた写像類群の表現からも、
その豊かな構造が垣間見れるということを紹介したい。

ウェブ図とは、経路積分で使われるファインマン図にいくつかの数学的な関係を入れたような
ある種のグラフであり、結び目の Kontsevich 不変量の構成で使われたものである。
Kontsevich は、数多くの先駆的な仕事をしているが、
その中で経路積分的に計算される量を組み合わせ論的な量であらわす（またはこの逆）ものがあり、
その一つが、コード図と KZ-方程式に関する反復積分とを用いて構成された Kontsevich 不変量である。
これは、Jones 多項式など、量子群から作られる量子不変量をほとんど全て含むと共に、
結び目の有限型不変量に関しても普遍的な、非常に強力な不変量である。
ウェブ図とは、このコード図を少しだけ一般化したものであり、
形式的には、量子群のもととなっている Lie 環の、グラフによる普遍的な表現と呼ぶべきものである。

さて、Kontsevich 不変量は、３次元多様体の不変量にも拡張できる。
これは、３次元多様体の量子不変量（Jones--Witten 不変量など）の摂動展開や、
有限型不変量について普遍的な性質をもち、普遍摂動不変量と呼ばれる。
さらに、境界付きの３次元多様体に対する普遍摂動不変量と、
その境界への写像類群の作用から、自然に、
ウェブ図のなす線形空間上への写像類群の作用が定義される。
そして、この作用の様子を見ると、
ウェブ図の持ついくつかの性質が写像類群の作用で保たれることがわかる。
これらの性質が、写像類群の幾何学的考察から導かれる性質とどのように対応するのかは
まだ良くわかっていないが、
このような性質には、ウェブ図のある種の頂点の数やオイラー数などがあり、
グラフの性質をあらわす数としては良く知られたものである。

この講演では、上で述べた構成について分かりやすく解説するとともに、
最後に、簡単な場合についてこの表現が具体的にどのようになっているかをウェブ図を用いて示す。
意外と複雑なことになっているなあと思ってもらえれば、
写像類群の豊かさを多少とも感じたことになるのではないかと思います。

\newpage

%======================================================================
% 講演６：中村博昭「ガロア・タイヒミュラーの夢」
%======================================================================
\begin{center}
{\large\bfseries ガロア・タイヒミュラーの夢}

\hspace*{\fill}{\large\bfseries 中村 博昭（都立大・理）}
\end{center}

\medskip

写像類群（＝タイヒミュラー群ともいう）をモジュライ空間の基本群とみる見方を、
グロタンディーク流の代数幾何的ガロア理論の観点から押し進めると、
自然な興味深い研究対象として現れてくるのが、ガロア・タイヒミュラー群——これは、
写像類群（のある種のコンパクト化）を絶対ガロア群と呼ばれる「クッション」で
包み込んだような巨大な位相群——である。

曲面の位相幾何的な切り貼りに応じて種数や境界成分の個数を変化させていく操作に応じて、
一連のガロア・タイヒミュラー群たちは「クッション部分」を共有しながら
様々な準同型写像で結ばれていく。
注意深く観察すると、ガロア・タイヒミュラー群の織りなす塔の本質的な情報の多くは、
小数の基本的なブロックとそれらの組み合わせ方に大きく依存している、
という構造が見えてくる（ガロア・タイヒミュラー塔の Lego 構造）。

こうした構造を利用して、ガロア群に関する整数論的な知識と
Johnson 準同型の様子に関する位相幾何的な情報とを結びつけたり、
また、絶対ガロア群のパラメータづけに関する方程式系を議論したり、
といった事が行なわれている昨今である。

\end{document}
